\documentclass[11pt]{article}

\usepackage[utf8]{inputenc}
\usepackage[T1]{fontenc}
\usepackage{amsmath,amssymb,amsthm}
\usepackage{booktabs}
\usepackage{array}
\usepackage{graphicx}
\usepackage[colorlinks=true,linkcolor=black,citecolor=blue,urlcolor=blue]{hyperref}
\usepackage[margin=1in]{geometry}

\newtheorem{definition}{Definition}
\newtheorem{proposition}{Proposition}
\newtheorem{remark}{Remark}

\newcommand{\sest}{\sigma}
\newcommand{\Mset}{\mathcal{M}}
\newcommand{\Cand}{\mathcal{C}}
\newcommand{\Rules}{\mathcal{R}}
\newcommand{\E}{\mathbb{E}}
\newcommand{\zero}{\mathsf{null}}

\title{\bf From Replayable Traces to Provenance-Aware Rule Fusion:\\
A Bounded Empirical Study of What RL Auditability Can and Cannot Identify}
\author{Working manuscript generated from repository evidence}
\date{September 2026}

\begin{document}
\maketitle

\begin{abstract}
\noindent
``Replayable traces,'' ``causally meaningful symbols,'' ``shared symbolic rules,'' and
``composable policies'' are distinct properties of a reinforcement-learning (RL) system,
yet they are frequently conflated under a single label of auditability. We treat them as
separate, individually testable claims and report a bounded study that measures each one
against executed artifacts. The work begins from a naive prototype whose grammar, replay,
and instrumentation claims did not survive scrutiny, and rebuilds the pipeline so that every
statement is anchored to a checkable quantity. On two discrete control tasks
(CartPole-v1 and Acrobot-v1), eight seeded agents per task share one frozen quantile state
symbolizer. Exact rule overlap between two CartPole agents is high (Jaccard $0.5833$, $77$ of
$132$ shared productions), yet a fresh raw-state behavioral test yields only $51.36\%$ action
agreement ($95\%$ episode-cluster interval $50.48\%$--$52.26\%$), a counterexample showing
that symbolic overlap is not behavioral consensus. A provenance-aware merger composes source
rule banks with explicit conflict arbitration, blind-spot detection, and fallback logging;
its hash-linked decisions resolve back to source rules and replay exactly against the
environment. CartPole fusion reaches $205.96$ mean return versus $163.65$ for the best single
actor; Acrobot fusion reaches $-443.37$ versus $-500$. Ablations expose strong, unexplained
sensitivity to arbitration, including an advantage for randomized rule selection, and a fitted
value-model baseline fails its quality checks. The evidence supports the \emph{existence} of
auditable trace and rule-composition mechanisms in these settings; generality, complete causal
explanation, and universal auditability of arbitrary RL remain unproven. A separate
designed-communication study reports a causal intervention on a learned symbol and is kept
evidence-disjoint from the rule-fusion conclusion.
\end{abstract}

\section{Introduction}

\paragraph{Scenario.}
A trained neural policy is often opaque: given a state it emits an action, but the internal
quantities that produced the choice are not directly legible. A growing body of work therefore
proposes to record RL execution as symbolic structure---replayable trajectories, induced
grammars, or state--action rule sets---so that a reviewer can inspect, reproduce, and reason
about a policy after training. The implicit promise is that a symbolic layer makes RL
``auditable.''

\paragraph{Limitations of existing approaches.}
That promise bundles several different questions, and the literature tends to answer one and
claim another. First, a lossless sequence grammar that reconstructs a \emph{seen} trace says
nothing about whether its productions generalize to unseen states or explain a decision.
Second, an overlap measure over symbolic rules is a statistic about syntax, not about behavior.
Third, a deterministic record can be internally consistent (a hash chain verifies) while the
transitions it encodes never happened; integrity is not fidelity. These are not hypothetical
worries: each one is instantiated as a failure we observed while building this system.

\paragraph{Problem and objective.}
Our objective is not to argue that an AI-native symbol system makes RL training fully
auditable. It is the narrower, falsifiable question: \emph{which specific quantities can an
induced symbolic layer plus an auditable execution record actually identify, and where does the
tempting inference from each quantity to a stronger one break?} We answer it by measuring each
quantity separately and reporting the counterevidence where an inference fails.

\paragraph{Challenges.}
Three challenges make this harder than it appears. (i) The naive intuition that ``compression
equals understanding'' fails at first contact: an early, readable serialization of the induced
grammar made the artifact $1.80\times$ \emph{larger} on disk than the raw trace it encoded, so
symbolic reduction did not deliver a storage benefit and could not be sold as compression.
(ii) Instrumentation is not free: an apparent difference between two training arms that should
have been identical turned out to be a hashing-definition error, not a real effect, and had to
be repaired before the equivalence claim was trustworthy. (iii) A large, precise-looking symbolic
core can still fail to predict behavior, so an overlap statistic must be paired with an
independent behavioral measurement.

\paragraph{Method.}
We address these by separating the pipeline into a \emph{generator} (a stochastic REINFORCE
agent that produces trajectories), a \emph{deterministic kernel} (state symbolization, rule
admission, conflict arbitration, and hash-chained logging), and a family of \emph{instruments}
(independent environment replay, a fresh-state behavioral-agreement test, and a provenance
resolver) that each target one estimand. Claims are then routed through an explicit budget:
positive, negative, diagnostic, and non-claim material are kept distinct.

\paragraph{Contributions.}
\begin{enumerate}
\item A provenance-aware rule-fusion mechanism whose decisions record source-rule attribution,
conflict arbitration, blind-spot detection, and fallback, and whose decision trace resolves back
to source rules and replays against the environment (Section~\ref{sec:fusion}).
\item A shared-symbolizer, multi-seed measurement of exact rule overlap paired with an
independent fresh-state behavioral-agreement test, yielding the counterexample that overlap does
not identify consensus (Section~\ref{sec:behavior}).
\item A failure-driven account of how a naive symbolic audit pipeline must be corrected to
support bounded claims, including an explicit pass/fail auditability taxonomy
(Sections~\ref{sec:pipeline} and~\ref{sec:instrument}).
\item Arbitration and value-model diagnostics that bound the policy-composition conclusions and
identify unresolved mechanisms (Section~\ref{sec:arbitration}).
\item A separate causal-intervention result for learned communication symbols in a designed task,
reported alongside---never merged into---the rule-fusion evidence (Section~\ref{sec:aim}).
\end{enumerate}

\section{From a Naive Prototype to a Validated Instrument}
\label{sec:pipeline}

This section is deliberately shallow-to-deep: it recounts the failures that shaped the final
measurement design, because those failures \emph{are} the reason the claims are bounded the way
they are.

\paragraph{Grammar conflation.}
The initial prototype treated two different objects as one. A SEQUITUR-style sequence
compressor, applied to a tokenized trajectory, produces a context-free grammar that can
\emph{reconstruct the string it was shown}; the source literature is explicit that this is a
coding of one observed sequence, not a policy. A separate state--action rule set
$\textit{state-symbol}\!\rightarrow\!\textit{action-symbol}$ is a candidate policy
\emph{generalization}. Merging them---and inferring rule provenance by substring tests rather
than by recorded occurrences---made the prototype unable to certify either replay or
generalization. The fix was architectural: the codec and the rule bank are induced, evaluated,
and reported separately throughout this paper.

\paragraph{Compression that was not compression.}
A natural first reading is that a smaller symbolic grammar means a smaller artifact. It did.
The readable serialization expanded a $61{,}295{,}430$-byte trace into a
$110{,}549{,}187$-byte grammar file ($1.80\times$), because pretty-printed string references cost
more than the tokens they named. After re-encoding with integer symbol identifiers and
general-purpose compression, the grammar archive was $78.2\%$ smaller than the raw trace but
still $8.4\%$ \emph{larger} than simply compressing the trace directly. The correct conclusion
is narrow: the grammar supplies exact reconstruction and explicit production structure, not a
best-in-class storage win. This becomes a permanent non-claim (Section~\ref{sec:claims}).

\paragraph{Integrity without fidelity.}
A hash chain over logged records detects tampering but does not prove the recorded transitions
occurred. We therefore added an \emph{independent} environment replay check: re-execute each
logged reset seed and action in the simulator and compare state, next state, reward, and the
terminated/truncated flags exactly. Replay is the identification-bearing test; the hash chain is
a diagnostic.

\paragraph{A hash bug that looked like a training difference.}
Two arms designed to train identically first appeared to differ. Inspection showed the
comparator had folded an arm label and a grammar-sidecar counter into the hash; the weights and
transitions were in fact equal. After hashing only the network state dictionary and defining a
chain-normalized core-update hash that excludes non-optimization sidecar fields, the arms matched
byte-for-byte. We record this correction rather than hiding it, because it is why the equivalence
claim is about \emph{optimizer-relevant} content only.

\section{Formal Setup}
\label{sec:setup}

\paragraph{Objects.}
An agent interacts with an episodic environment producing a state $s_t$, an action $a_t$, a
reward $r_t$, a next state $s_{t+1}$, and terminal and truncated flags $\mathsf{term}_t$ and
$\mathsf{trunc}_t$. The observation $s_t$ is a real vector;
the action $a_t$ is drawn from a finite set $\mathcal{A}=\{0,\dots,|\mathcal{A}|-1\}$. A policy
is a stochastic map $\pi_\theta(a\mid s)$ trained by REINFORCE with an entropy-free sampled
objective and a normalized discounted return with $\gamma=0.99$.

\begin{definition}[State symbolization]
\label{def:symbolizer}
A \emph{state symbolizer} $\sest$ maps a real state to a tuple of bin indices
\begin{equation}
\sest(s) \;=\; \big(\,\mathrm{bucketize}_{d}\!\big(s^{(d)};\,\theta_{d}\big)\,\big)_{d=1}^{D},
\label{eq:symbolize}
\end{equation}
where each $\theta_{d}$ holds the interior empirical-quantile edges of dimension $d$ over $B$
bins. The edges are fitted once from pooled random-policy calibration rollouts and then
\emph{frozen}, so that all training seeds and arms within a task share the same symbol alphabet.
\end{definition}

\begin{definition}[Admitted state--action rule]
\label{def:rule}
Let $N(c,a)$ be the count of observed transitions whose symbol is $c=\sest(s)$ and action is $a$,
and $N(c)=\sum_{a}N(c,a)$. The induced action for condition $c$ is
$a^{\star}(c)=\arg\max_{a}N(c,a)$ (ties to the smallest index), with support
$\kappa_{s}(c)=N(c,a^{\star})$ and confidence
\begin{equation}
\kappa(c) \;=\; \frac{N(c,a^{\star}(c))}{N(c)},
\qquad
\text{rule } c\!\rightarrow\!a^{\star}(c) \text{ is admitted iff }
\kappa_{s}(c)\ge \kappa_{s}^{\min}\ \wedge\ \kappa(c)\ge \tau .
\label{eq:rule}
\end{equation}
The admitted rules form a rule bank $\Rules$. Each rule stores the multiset of source step
identifiers that produced it, and the digest of that multiset.
\end{definition}

Equation~\eqref{eq:rule} uses a single normalization for confidence; support and confidence are
reported separately from coverage so that a precise-but-rare rule is never described as broadly
applicable.

\section{Mechanisms}
\label{sec:mechanisms}

\subsection{A grammar-constrained policy}
Given an active bank $\Rules$ from Definition~\ref{def:rule}, a constrained policy masks the
action set to the rule's action when the current symbol is covered and leaves it unconstrained
otherwise:
\begin{equation}
\Mset(c)=\begin{cases}\{a^{\star}(c)\}, & c\in\Rules,\\[2pt]\mathcal{A}, & c\notin\Rules,\end{cases}
\qquad
\pi_{\mathrm{c}}(a\mid s)=\frac{\exp z_{a}(s)\,\mathbf{1}[a\in \Mset(\sest(s))]}
{\sum_{a'\in\Mset(\sest(s))}\exp z_{a'}(s)} .
\label{eq:constrained}
\end{equation}
The mask is a \emph{state-conditioned} restriction with an explicit unconstrained fallback, which
corrects an earlier global action-pruning design that could permanently remove an action merely
because it was absent from frequent rules.

\subsection{Hash-chained records and replay validity}
\label{sec:replaymech}
Each logged event $e_t$ commits to its predecessor:
\begin{equation}
h_t=\mathrm{SHA256}\!\big(h_{t-1}\,\|\;\mathrm{canon}(e_t)\big),\qquad h_{-1}=\zero,
\label{eq:chain}
\end{equation}
so modifying any $e_t$ breaks every later digest. Integrity of Equation~\eqref{eq:chain} is a
diagnostic. The identification-bearing predicate is \emph{environment replay validity}: a trace
is replay-valid iff, re-executing recorded reset seeds and actions in the simulator, every step's
state, next state, reward, and both terminal flags match the record exactly (compared as
32-bit floating-point quantities and Boolean flags). A hash chain can be internally consistent
while replay fails; only the latter supports a transition-fidelity statement.

\subsection{Provenance-aware rule fusion}
\label{sec:fusion}
At a state $s$ with symbol $c=\sest(s)$ and $K$ source rule banks $\Rules_1,\dots,\Rules_K$, the
merger forms the candidate set
\begin{equation}
\Cand(c)=\{\,i : c\in\Rules_i\,\},\quad
\text{covered}\iff|\Cand(c)|\ge1,\quad
\text{conflict}\iff\big|\{a_i:i\in\Cand(c)\}\big|>1 .
\label{eq:fusion}
\end{equation}
Covered non-conflicts and agreements emit the common action; conflicts are resolved by the
declared priority order $(\text{confidence},\text{support},\text{mean reward},-\text{seed})$;
uncovered states take an actor fallback. Each decision records both candidate rule identifiers,
their full statistics and source-step digests, and the selected source or fallback, and enters a
chain of the form of Equation~\eqref{eq:chain}. A validator resolves every selected reference
back to its source bank, re-checks the recorded statistics and emitted action, verifies the chain,
and replays the resulting actions in the environment. Blind-spot selection uses training-time
returns only, so no evaluation return leaks into the fallback choice.

\subsection{Behavioral agreement and symbolic overlap}
\label{sec:behavior}
For two frozen actors $\pi_1,\pi_2$ and a fresh held-out state distribution $\mathcal{D}$ sampled
equally from each actor's own occupancy, define
\begin{equation}
A \;=\; \Pr_{s\sim\mathcal{D}}\!\big[\arg\max_a \pi_1(a\mid s)=\arg\max_a \pi_2(a\mid s)\big],
\label{eq:agreement}
\end{equation}
estimated on $10{,}000$ fresh raw states and bracketed by an episode-cluster bootstrap. The
symbolic-overlap statistic is exact-set Jaccard over $(\text{state-symbol},\text{action-symbol})$
productions of the two banks,
\begin{equation}
J \;=\; \frac{|\Rules_1\cap \Rules_2|}{|\Rules_1\cup \Rules_2|}.
\label{eq:jaccard}
\end{equation}
Equations~\eqref{eq:agreement} and~\eqref{eq:jaccard} are measured on the same agents precisely so
that one cannot be substituted for the other.

\subsection{Trajectory coding length}
The lossless codec induces pair-substitution productions and reports a declared two-part
description length
\begin{equation}
L_{\mathrm{mdl}} \;=\; \underbrace{L(\Rules_{\mathrm{cfg}})}_{\text{grammar}} \;+\;
\underbrace{L(\tau\mid \Rules_{\mathrm{cfg}})}_{\text{start sequence}},
\label{eq:mdl}
\end{equation}
with Elias-gamma-coded counts and literal terminal dictionary, and fixed-width symbol references;
provenance bytes are reported separately. Equation~\eqref{eq:mdl} is a description-length score
under a stated convention, not an optimization objective.

\subsection{Causal symbol intervention (separate sub-study)}
\label{sec:aim}
In an independent two-agent private-target coordination game, a receiver maps a state and a
discrete sender message to an action. The matched-support intervention contrasts the receiver's
action distribution under the observed message $m$ against a counterfactual message $m'$ that
keeps the information-bearing symbol's context fixed:
\begin{equation}
\Delta \;=\; \E_{s\sim\mathcal{S}}\!\left[\,\sum_{a}\pi_R(a\mid s,m) - \pi_R(a\mid s,m')\,\right],
\label{eq:interv}
\end{equation}
evaluated only on the common support of the two conditioned state sets. Equation~\eqref{eq:interv}
is computed for the AIM sub-study alone and is not merged into the rule-fusion estimands.

\section{The Instrument and What It Identifies}
\label{sec:instrument}

We separate a set of quantities that a single ``auditability score'' would collapse
(Table~\ref{tab:taxonomy}). Trace integrity (Equation~\eqref{eq:chain}) is tamper-\emph{evidence}:
a counterfactual in which a modified record still passes is unreachable by construction and so is
not identification-bearing. Environment replay is identification-bearing but only for the recorded
executions under the pinned software and hardware configuration. Lossless coding
(Equation~\eqref{eq:mdl}) is a reconstruction statement. Rule coverage and agreement measure how
often a rule applies and how often it matches a source actor. Behavioral agreement
(Equation~\eqref{eq:agreement}) is an independent action-level estimand. Composition quality
(Equation~\eqref{eq:fusion}) is a measured return under declared rules. Value-model reliability is
a separate gate the fitted critics in this study do not pass.

\begin{table}[t]
\centering\small
\begin{tabular}{@{}p{0.26\linewidth}p{0.30\linewidth}p{0.34\linewidth}@{}}
\toprule
Quantity & Estimand / predicate & What it cannot establish \\
\midrule
Trace integrity & Chain~\eqref{eq:chain} verifies & transitions occurred; logger honesty \\
Environment replay & exact transition match & cross-hardware equivalence \\
Lossless coding & reconstruction~\eqref{eq:mdl} & policy explanation; storage win \\
Rule coverage/agreement & support/confidence~\eqref{eq:rule} & generalization at low coverage \\
Behavioral agreement & $A$ in~\eqref{eq:agreement} & optimality; causal rationale \\
Symbolic overlap & $J$ in~\eqref{eq:jaccard} & behavioral consensus \\
Composition quality & return of~\eqref{eq:fusion} & optimal arbitration \\
\bottomrule
\end{tabular}
\caption{Auditability disaggregated into separately measured quantities. Each row's right column is
a boundary enforced by a Phase-0 negative or diagnostic result.}
\label{tab:taxonomy}
\end{table}

\paragraph{Transport versus generation.}
For the AIM sub-study, the intervention changes the message content presented to the receiver, not
its transport; it is therefore classified as a generation-relevant change. For the RL kernel, the
merger neither samples nor learns; it is a deterministic decision procedure. A change of record
serialization is never treated as a change of policy.

\section{Experimental Design and Registered Predictions}
\label{sec:design}

Eight seeds $\{11,29,43,71,101,149,211,307\}$ train REINFORCE agents on CartPole-v1 and
Acrobot-v1 for $300$ episodes per seed and arm, under one frozen shared quantile symbolizer per
task ($B=4$ bins, calibration rollouts pooled across seeds). Rule admission uses support
$\kappa_s^{\min}=8$ and confidence $\tau=0.70$. A baseline arm and an audit-only observer arm are
checked for exact optimizer-level equivalence. Composition uses $100$ common reset seeds per
policy; Acrobot arbiter ablations reuse those seeds, and randomized rule selection repeats over
five independent selector streams. A separate five-stream uniform-random-action control isolates
unconditioned action noise from state-conditioned rule availability. Stochastic return summaries
use the across-episode standard deviation with a percentile bootstrap over episodes; state-level
agreement treats sampled states as the measurement unit but clusters the bootstrap by episode.

Four hypotheses were recorded before the corresponding measurements:
\begin{itemize}
\item[\textbf{H1}] A grammar induced from training symbols reconstructs those trajectories exactly,
but held-out episodes are not ``replayed'' merely because their tokens are in the alphabet.
\item[\textbf{H2}] A learned state--action rule bank predicts an independent holdout above an
action-frequency baseline, with coverage, agreement, and confidence reported separately.
\item[\textbf{H3}] Grammar constraints applied only in well-supported states do not harm return
when paired over seeds, with an explicit low-coverage fallback.
\item[\textbf{H4}] Reconstructing a sampled trace alone is insufficient for whole-training
auditability; configuration, transitions, update boundaries, and checkpoints must be recorded and
re-derived.
\end{itemize}
Each prediction is scored in Section~\ref{sec:results} as supported, partially supported, or
contradicted. Because the measurements were exploratory in the final shared-symbolizer batch, the
H-statuses are reconciliations against pre-specified statements rather than a formal
preregistration.

\section{Results}
\label{sec:results}

\subsection{Overlap versus behavior (H2)}
\label{sec:behaviorres}
\emph{Prediction.} H2 anticipated that a large shared rule core would coincide with faithful
cross-seed behavior. \emph{Result.} It does not. Across eight agents, mean pairwise exact overlap
is $J=0.5572$ on CartPole and $0.0956$ on Acrobot. For seeds 11 and 29, CartPole has $77$ shared
productions out of a $132$ union ($J=0.5833$) with no action conflicts on shared conditions;
Acrobot has $21$ of $423$ ($J=0.0496$) with $77$ of $98$ shared conditions in action conflict.
On $10{,}000$ fresh CartPole raw states, the seed-11 and seed-29 actors agree on $51.38\%$ of
argmax actions; the episode-cluster balanced estimate is $51.36\%$ ($95\%$ interval
$50.48\%$--$52.26\%$), near chance for a two-action task. This is a central negative result and
\textbf{contradicts} H2: a substantial exact intersection does not imply behavioral consensus,
and $J$ alone is disqualified as evidence of a shared policy core.

\subsection{Trace and provenance auditability (H1, H4)}
For the recorded artifacts, chain, update-ledger, provenance-resolution, and environment-replay
checks pass under the audited setup. The eight-seed records span $1{,}725{,}655$ CartPole and
$2{,}544{,}513$ Acrobot trace rows across three arms; replay covers $1{,}122{,}403$ and
$1{,}742{,}570$ transitions respectively. Every logged prefix round-trips exactly
(Equation~\eqref{eq:mdl}), and a tamper probe that edits a single reward is rejected at the first
line. These support \textbf{H1} and a bounded \textbf{H4}: exact replay is evidence about recorded
executions and specified software/hardware conditions; it does not establish that the logger
captured every training cause, prove logger honesty, or guarantee cross-hardware equivalence.

\subsection{Composition returns and arbitration sensitivity (H3)}
\label{sec:arbitration}
\emph{Prediction.} H3 anticipated that well-supported constraints with a fallback would not harm
return. \emph{Result.} The online majority constraint is unstable: in the two-seed CartPole pilot,
the constrained-arm return changed by $-378.7$ for seed 11 but $+31.5$ for seed 29 (mean
$-173.6$, disagreeing signs); lowering the threshold to $0.70$ gave $-330.4$ and $+17.0$. This
\textbf{contradicts} H3 as a performance claim and motivates reporting composition returns as
descriptive only.

Table~\ref{tab:returns} gives the shared-symbolizer composition returns. CartPole fusion reaches
$205.96$ versus $163.65$ for the best single actor (paired difference $+42.31$, $95\%$ bootstrap
$[23.98,61.93]$), but some policies hit a $500$-step ceiling that masks differences. Acrobot
fusion reaches $-443.37$ versus $-500$ (paired $+56.63$, $95\%$ $[39.32,75.15]$) at $95.70\%$
rule coverage, $89.24\%$ conflict, and $4.30\%$ blind spots; $39$ of $100$ episodes terminate
early, though the median remains $-500$.

\begin{table}[t]
\centering\small
\begin{tabular}{@{}llrrr@{}}
\toprule
Task & Policy & Mean return & Paired $\Delta$ vs best & $95\%$ CI \\
\midrule
CartPole & best single actor   & $163.65$ & --- & --- \\
CartPole & actor-logit ensemble & $500.00$ & --- & --- \\
CartPole & grammar fusion       & $205.96$ & $+42.31$ & $[23.98,61.93]$ \\
Acrobot  & best single actor    & $-500.00$ & --- & --- \\
Acrobot  & actor-logit ensemble & $-500.00$ & --- & --- \\
Acrobot  & grammar fusion       & $-443.37$ & $+56.63$ & $[39.32,75.15]$ \\
\bottomrule
\end{tabular}
\caption{Composition returns over $100$ matched reset seeds. CartPole is ceiling-saturated, so the
ensemble comparison is not a controlled superiority test.}
\label{tab:returns}
\end{table}

Arbiter ablations (Table~\ref{tab:arb}) show strong, unexplained sensitivity. The declared
order on confidence, support, and mean reward scores $-443.37$; confidence-only $-444.49$ (paired $-1.12$,
$95\%$ $[-3.16,+0.06]$); mean-reward-only, support-first, and conflict deferral to the ensemble
each score $-500$. Randomly choosing among available source rules scores $-187.02$ on one stream
and $-194.10$ to $-206.17$ across five (mean $-198.82$), while a five-stream uniform-random-action
control averages $-498.71$. The mean-reward key is near-inert here ($1575$ of $1694$ Acrobot rules
have mean reward exactly $-1.0$). The random-action control separates state-conditioned rule
availability from unconditioned noise, but the mechanism of the randomized-selection advantage is
\emph{not} identified, and the tested confidence-ranked arbiter is not shown to be a strong
conflict resolver.

\begin{table}[t]
\centering\small
\begin{tabular}{@{}lrr@{}}
\toprule
Acrobot arbiter variant & Mean return & Paired $\Delta$ vs declared \\
\midrule
confidence/support/reward (declared) & $-443.37$ & --- \\
confidence-only                      & $-444.49$ & $-1.12$ \\
mean-reward-only                     & $-500.00$ & $-56.63$ \\
support-first                        & $-500.00$ & $-56.63$ \\
defer conflicts to ensemble          & $-500.00$ & $-56.63$ \\
random rule selection (5 streams)    & $-198.82$ & $+244.55$ \\
uniform random action (5 streams)    & $-498.71$ & $-55.34$ \\
\bottomrule
\end{tabular}
\caption{Matched-seed arbitration ablations. Random rule selection beats the declared arbiter;
random action does not, so the effect tracks state-conditioned availability, not action noise.}
\label{tab:arb}
\end{table}

\subsection{Value-model diagnostics (kept separate)}
The fitted value baseline fails its quality gates and is reported apart from the fusion result. On
CartPole it achieves mean return $9.34$ (versus $163.65$ for the best actor) and matches the
actor-logit ensemble on $35.27\%$ of $31{,}563$ held-out states. On Acrobot it scores $-499.57$
and selects action $0$ on all $50{,}000$ diagnostic states. On $200$ Acrobot states the seed-11 and
seed-29 critics correlate at $0.652$ in scalar value while their greedy actions diverge; fusion
matches the seed-11 argmax on $64.5\%$, the seed-29 argmax and multi-critic improvement on $0\%$,
and the best first action under a fixed ensemble continuation on $22.5\%$ ($95\%$ interval
$16.5\%$--$28.5\%$). A scalar-value correlation is not action agreement; these critics are an
unvalidated comparator, not evidence against the grammar-fusion return.

\subsection{Causal symbol intervention (separate sub-study)}
\label{sec:aimres}
\emph{Prediction.} None registered for the rule-fusion estimator. \emph{Result.} Across six seeds
and $600{,}000$ logged decisions in the designed coordination game, sender-sampled receiver
accuracy averages $0.998923$, versus $0.503292$ for shuffled messages and $0.500$ for a no-message
control. The matched-support intervention~\eqref{eq:interv} on the information-bearing symbol
changes receiver action probability by $0.998366$ on average (per-seed $0.996986$--$0.999685$;
mean common-support overlap $0.94308$), and the active token position varies across seeds. This
supports a causal role for a learned symbol \emph{in the tested game}; it does not provide a
general semantic decoding of the language, and its executed artifacts live outside the compact
review package, so it is kept evidence-disjoint from Sections~\ref{sec:behaviorres}
and~\ref{sec:arbitration}.

\subsection{Pass/fail summary}
\begin{table}[t]
\centering\small
\begin{tabular}{@{}lllll@{}}
\toprule
Criterion & Estimate & Threshold & Verdict & Layer \\
\midrule
Replay match (CartPole) & $1{,}122{,}403$ transitions & exact & pass & identification \\
Replay match (Acrobot)  & $1{,}742{,}570$ transitions & exact & pass & identification \\
Baseline/observer equivalence & byte-identical core & exact & pass & health \\
Lossless reconstruction & ratio $0.513$--$0.575$ & round-trip & pass & diagnostic \\
Storage vs.\ generic compression & $+8.4\%$ & smaller & fail & diagnostic \\
Behavioral agreement $A$ & $51.36\%$ & $\gg$ chance & fail & identification \\
Symbolic overlap $J$ (CartPole) & $0.5833$ & descriptive & reported & descriptive \\
Value-model gate & action $0$ everywhere & informed & fail & diagnostic \\
\bottomrule
\end{tabular}
\caption{Consolidated verdicts. The single identification-bearing failure (behavioral agreement)
and the identification-bearing passes (replay) are separated from health and diagnostic rows.}
\label{tab:passfail}
\end{table}

Identification-bearing outcomes (replay fidelity; the overlap-versus-behavior counterexample) are
reported apart from operational-health rows (arm equivalence), calibration/diagnostic rows (hash
chain, lossless coding), and the failed value gate.

\section{Related Work}
\label{sec:related}

\paragraph{Direction: interpretable and programmatic policies.}
Programmatically interpretable RL represents policies in a domain-specific language and reasons
about them symbolically. Our setting differs: rules are induced from logged traces, so provenance
and predictive generalization must each be earned, not assumed. That requirement is exactly where
our negative result bites---syntax overlap does not grant behavioral faithfulness.

\paragraph{Methodology: sequence structure and policy composition.}
Linear-time grammar induction identifies hierarchical structure in sequences but, by the original
authors' own caveat, codes a single observed string rather than a generalized policy. We keep that
distinction operational (Definition~\ref{def:rule} versus Equation~\eqref{eq:mdl}). Generalized
policy improvement bounds a composed policy by the maximum per-state advantage of its sources;
successor-feature composition likewise presumes reasonably fitted value functions. Our fitted
critics violate that premise (Section~\ref{sec:arbitration} diagnostics), so we invoke the
composition \emph{idea} only as a diagnostic baseline, not as a bound we establish.

\paragraph{Measurement of emergent communication and reproducibility.}
Pitfalls in measuring emergent communication motivate treating a designed-game causal test and a
state--action rule test as separate estimands. Deterministic-implementation analysis for deep RL
separates deterministic code paths from same-condition replicability, which is why our replay
claim is bounded to the recorded host and versions rather than generalized. Where these lines of
work and ours could conflict---for instance on whether symbolic overlap should predict behavior---
our evidence is a single-task counterexample, so we mark it as a bounded observation, not a
refutation of a broad claim.

\bibliographystyle{plain}
% The references are embedded below to keep the manuscript self-contained.
\begin{thebibliography}{9}
\bibitem{verma2018pirl}
Verma, Murali, Singh, Kohli, and Chaudhuri. Programmatically interpretable reinforcement learning.
\emph{ICML}, PMLR 80, 2018.
\bibitem{nevill1997sequitur}
Nevill-Manning and Witten. Identifying hierarchical structure in sequences: a linear-time algorithm.
\emph{Journal of Artificial Intelligence Research}, 7:67--82, 1997.
\bibitem{thakoor2022gpi}
Thakoor, Rowland, Mavrin, Borsa, Dabney, Munos, and Barreto. Generalised policy improvement with
geometric policy composition. \emph{ICML}, PMLR 162, 2022.
\bibitem{barreto2017sf}
Barreto, Dabney, Munos, Hunt, Schaul, van Hasselt, and Silver. Successor features for transfer in
reinforcement learning. \emph{NeurIPS} 30, 2017.
\bibitem{lowe2019pitfalls}
Lowe, Foerster, Boureau, Pineau, and Dauphin. On the pitfalls of measuring emergent communication.
\emph{AAMAS}, 2019. arXiv:1903.05168.
\bibitem{nagarajan2018det}
Nagarajan, Warnell, and Stone. Deterministic implementations for reproducibility in deep
reinforcement learning. arXiv:1809.05676, 2018.
\end{thebibliography}

\section{Discussion}
\label{sec:discussion}

\paragraph{Why the mechanism can be trusted as far as it is.}
The audit capability rests on two equations that are independently checkable: the chain
commitment~\eqref{eq:chain} makes tampering detectable, and the replay predicate makes transition
fidelity verifiable in the environment. The merger's decisions are auditable because each carries
its source references, which resolve back to the inducing statistics. None of this depends on the
grammar being an explanation; it depends on the record being reconstructible and re-executable.

\paragraph{Expected versus unexpected outcomes.}
Expected: replay fidelity, high rule coverage, and the value-model fragility. Unexpected: (i) a
large exact overlap coexisting with near-chance behavior ($51.36\%$ at $J=0.5833$), and (ii)
randomized rule selection beating the confidence-ranked arbiter by well over two hundred return
points. For (i) the alternative explanation is that many distinct rule sets induce the same, and
different, argmax behaviors once states are quantized differently; our data cannot separate this
from genuine divergence. For (ii) a plausible alternative is that the advantage comes from the
stochastic mixture escaping the $-500$ truncation trap, since candidate availability is
state-conditioned---but the uniform-action control ($-498.71$) shows it is not unconditioned noise
alone. We leave the mechanism unidentified.

\paragraph{Implications.}
Theoretically, the results caution against treating a syntax statistic as a semantic or behavioral
one. Practically, they suggest that provenance-carrying composition is inspectable without being
optimal. Neither implication claims downstream deployment quality we did not measure.

\paragraph{Steelman.}
The strongest objection is that the headline fusion gain is an artifact of a saturated comparison:
both best-single and ensemble sit at the Acrobot $-500$ cap, so almost anything that varies actions
with the state will appear better, and the randomized-selection result makes the ``intelligent''
arbiter look incidental. We take this seriously: we do not claim the arbiter is a good conflict
resolver, we report that random beats it, we disclose the ceiling on CartPole, and we keep the
improvement language descriptive of these policies and reset seeds. What survives the objection is
smaller but real: the composition trace is auditable and replayable, and its failure mode
(overlap $\neq$ behavior) is measured directly rather than inferred.

\paragraph{Replication (bounded to the repository).}
Independent replication requires the full multi-gigabyte trace and checkpoint tree; the compact
review package supplies the summaries, overlap tables, arbitration ablations, the random-action
control, Experiment B, and the value diagnostics with recorded digests. Minimum conditions are:
the pinned interpreter, tensor, and environment-library versions; deterministic single-thread CPU
settings; the shared calibration quantiles; and re-derivation of every reported digest. Replaying
a recorded output is distinguishable from generating a fresh one; only the latter tests transfer.

\section{Limitations and Future Work}
\label{sec:claims}
\textbf{Non-claims.} We do not claim: general or complete auditability of arbitrary RL training;
that symbolic overlap entails behavioral consensus; that induced grammars are human-readable
explanations of a neural policy; a storage-compression advantage over generic compression; optimal
or mechanism-explained arbitration; validated value-model composition; or a natural-language
semantics for the communication symbols. \textbf{Boundaries.} Evidence covers two benchmark tasks,
eight seeds per task, one shared quantizer per task, and one software/hardware configuration;
CartPole saturates and many Acrobot episodes remain at $-500$; fresh-state behavior is measured for
the seed-11/29 CartPole pair, not every pair or task. \textbf{Future work} follows directly from
the unresolved mechanisms: decompose the random-arbitration advantage by source rule, action
frequency, and state occupancy; test return-weighted rule-bank induction under a pre-specified
normalization and held-out evaluation; use the native sequence symbols as the control
representation; and evaluate behavior on non-saturated tasks and across hardware.

\section{Conclusion}
In the evaluated prototype, executed evidence supports replayable trace inspection,
provenance-aware rule composition, explicit conflict and blind-spot reporting, and a causal symbol
effect in a designed communication game. The same evidence shows that grammar overlap did not
predict actor agreement, that arbitration behavior remains unexplained, and that a fitted value
baseline failed its checks. We therefore claim bounded existence, not universality: these
mechanisms make specified RL artifacts more auditable under tested conditions, while general
auditability and complete causal explanation remain open.

\end{document}
